\begin{proof}
Write $C=C(m,2/3)$.  The established Hardy bridge gives
$E_m=1/C$.  Every minimizing model space has exact dimension $m$:
otherwise adjoining a reproducing-kernel direction with nonzero
overlap with the residual strictly improves its fit.  Write its
Blaschke product as $\Theta=D^*/D$, with $D$ an outer polynomial of
degree at most $m$, and set $b=P_{\le m}\Theta$.
The coefficient identity and H\"older inequality used in the bridge
are
\begin{equation}
 \Rnorm D_2^2=[bD^2]_m
 \le C\Rnorm{bD^2}_{2/3}
 \le C\Rnorm b_2\Rnorm D_2^2.
 \label{I:R:eq:classification-Holder-chain}
\end{equation}
At a minimizing space $\Rnorm b_2=E_m=1/C$, so equality holds
throughout.  In particular normalized $bD^2$ is a coefficient
extremal.  It is outer by the stated hypothesis, hence $b$ is outer
as well.  Equality in H\"older gives
$|b|/\Rnorm b_2=|D|/\Rnorm D_2$ almost everywhere on the circle.
Uniqueness of an outer function with its boundary modulus therefore
gives $b=cD$ for a nonzero scalar $c$.  It follows that the normalized
cube of $D$ is a coefficient extremal: $D$ is a scalar multiple of
some $P\in\mathcal P_m$.  Conversely the stationarity identity
$P^2=C P^*+O(z^{m+1})$ in the upper half of the Hardy bridge makes
$P^*/P$ a minimizing Blaschke product.  This proves the asserted
classification.

For a lossless observable boundary tuple $(D,l,x)$, its observation
map is an isometry onto the minimizing model space and
$O^*e_m=(lD^m)^*=\Gamma_m x$.  Thus the model space fixes the tuple
up to a unitary state change and fixes the unit cyclic vector by this
display.  A rotation of the scalar variable only produces the
corresponding phase changes in the tuple.
\end{proof}
